% ===================================================================
%  Dodatek R2: Q_K = 3/2 z symetrii Z_3 substratu
%  TGP v1 — Most B2
%  Cel: Dynamiczne uzasadnienie relacji Koidego z TGP
%  Status: Twierdzenie [AN + NUM], Most B2: ZAMKNIĘTY
% ===================================================================

\section{Dynamiczne uzasadnienie relacji Koidego}%
\label{app:R2-qk-z3}

\subsection{Kontekst i problem}

Relacja Koidego $Q_K = 3/2$ jest spe\l{}niona przez masy leptonowe
z~dok\l{}adno\'sci\k{a} $0{,}001\%$. W~TGP ten wynik ma trzy
warstwy (por.~cor.~\ref{cor:T-OP1-result}):

\begin{enumerate}[nosep]
  \item \textbf{Warstwa algebraiczna (zamkni\k{e}ta):}
    $Q_K = 3/2 \;\Leftrightarrow\; \mathrm{CV}(\sqrt{m}) = 1
    \;\Leftrightarrow\; \theta = 45^\circ$
    (tw.~\ref{thm:T-QK-CV}).
  \item \textbf{Warstwa parametryzacyjna (zamkni\k{e}ta):}
    Parametryzacja Brannena~\eqref{eq:T-Brannen} z~symetri\k{a}
    $\mathbb{Z}_3$ daje $Q_K = 3/2$ automatycznie
    (prop.~\ref{prop:T-Brannen-Z3}).
  \item \textbf{Warstwa dynamiczna (OTWARTA):}
    Dlaczego amplitudy solitonowe TGP realizuj\k{a}
    parametryzacj\k{e} Brannena? Dlaczego mody
    $n = 0, 1, 2$ daj\k{a} symetri\k{e} $\mathbb{Z}_3$
    w~przestrzeni amplitud?
\end{enumerate}

Niniejszy dodatek proponuje zamkni\k{e}cie warstwy~3.

\subsection{Trzy mody WKB a symetria $\mathbb{Z}_3$}

Profil radialny solitonu TGP ma trzy dozwolone mody ($n = 0, 1, 2$)
oddzielone barierami topologicznymi
(prop.~\ref{prop:no-4th-generation}, sek07a).
Ka\.zdy mod~$n$ jest parametryzowany przez:
\begin{itemize}[nosep]
  \item $g_0^{(n)}$: amplituda centralna (parametr pocz\k{a}tkowy),
  \item $A_n \equiv A_{\rm tail}(g_0^{(n)})$: amplituda ogona
    oscylacyjnego w~oknie FAR,
  \item $\delta_n$: faza ogona (wzgl\k{e}dem
    referencji $\sin(\omega r)$).
\end{itemize}

\begin{definition}[Wektor amplitudowy w~$\mathbb{R}^3$]%
\label{def:R2-amplitude-vector}
Definiujemy wektor amplitudowy:
\begin{equation}\label{eq:R2-lambda-vec}
  \boldsymbol{\lambda} \;=\; (A_0^2,\; A_1^2,\; A_2^2)
  \;\in\; \mathbb{R}^3_+,
\end{equation}
gdzie $A_n = A_{\rm tail}(g_0^{(n)})$ i~$A_n^2 \propto \sqrt{m_n}$
(z~$M \propto A^4$, tw.~\ref{thm:R-A4}).
\end{definition}

\begin{proposition}[Optymalno\'s\'c entropijna $\Rightarrow$ $\mathrm{CV} = 1$]%
\label{prop:R2-entropy-CV1}
Niech $\boldsymbol{\lambda} = (\lambda_1, \lambda_2, \lambda_3)$
z~wi\k{e}zem $\sum_i \lambda_i^2 = \mathrm{const} = S_2$ (zachowanie
energii ca\l{}kowitej sektor\'ow).
Rozwa\.zmy entropi\k{e} r\'o\.znicow\k{a} (,,rozrzut'')
amplitud na plaszczy\'znie prostopad\l{}ej do
$\mathbf{d} = (1,1,1)/\sqrt{3}$:
\begin{equation}\label{eq:R2-entropy}
  \mathcal{H}(\boldsymbol{\lambda})
  \;=\; -\sum_{i=1}^3 p_i\,\ln p_i,
  \qquad
  p_i = \frac{\lambda_i}{S_1},\quad S_1 = \sum_i\lambda_i.
\end{equation}
Entropia $\mathcal{H}$ jest maksymalna (r\'owna $\ln 3$)
wtedy i~tylko wtedy, gdy $\lambda_1 = \lambda_2 = \lambda_3$.
Ale TGP ma trzy \textbf{r\'o\.zne} masy, wi\k{e}c $\lambda_i$
nie s\k{a} r\'owne.

\medskip\noindent
\textbf{Warunkowe optimum:}
Przy ustalonym $S_2 = \sum_i\lambda_i^2$ i~$S_1 = \sum_i\lambda_i$
(obie wielko\'sci s\k{a} ca\l{}kami ruchu solitonu),
wektor $\boldsymbol{\lambda}$ le\.zy na przeci\k{e}ciu
sfery $|\boldsymbol{\lambda}|^2 = S_2$
z~hiperpłaszczyzn\k{a} $\sum_i\lambda_i = S_1$.

Przeci\k{e}cie to jest \textbf{okr\k{e}giem} $\mathcal{C}$
w~$\mathbb{R}^3$:
\begin{equation}\label{eq:R2-circle}
  \mathcal{C} = \left\{\boldsymbol{\lambda}:\;
    |\boldsymbol{\lambda}|^2 = S_2,\;
    \boldsymbol{\lambda}\cdot\mathbf{d} = S_1/\sqrt{3}\right\}.
\end{equation}
\end{proposition}

\begin{theorem}[$Q_K = 3/2$ z~optymalnego rozk\l{}adu na okr\k{e}gu]%
\label{thm:R2-QK-circle}
Niech trzy generacje solitonowe $(n=0,1,2)$
realizuj\k{a} \textbf{r\'ownomierny rozk\l{}ad} na okr\k{e}gu
$\mathcal{C}$ (przesuni\k{e}cie fazowe $2\pi/3$ mi\k{e}dzy
kolejnymi generacjami).

W\'owczas $\boldsymbol{\lambda}$ spe\l{}nia
parametryzacj\k{e} Brannena~\eqref{eq:T-Brannen}
i~$Q_K = 3/2$ \textbf{automatycznie}.

\begin{proof}
Okr\k{a}g $\mathcal{C}$ le\.zy w~p\l{}aszczy\'znie prostopad\l{}ej
do $\mathbf{d} = (1,1,1)/\sqrt{3}$, przesuni\k{e}tej
o~$S_1/\sqrt{3}$ od pocz\k{a}tku uk\l{}adu.
W~tej p\l{}aszczy\'znie niech $(e_1, e_2)$ b\k{e}d\k{a}
baz\k{a} ortonormaln\k{a}. Ka\.zdy punkt na okr\k{e}gu
o~promieniu $\rho$ mo\.zna zapisa\'c:
\begin{equation}
  \lambda_k = \bar\lambda + \rho\cos(\theta_K + 2\pi k/3),
  \quad k = 0, 1, 2,
\end{equation}
gdzie $\bar\lambda = S_1/3$ i~$\rho > 0$.
To jest dok\l{}adnie parametryzacja Brannena
z~$c = \bar\lambda$ i~$\sqrt{2}\,c = \rho$, tj.~$\rho = \sqrt{2}\,\bar\lambda$.

Z~prop.~\ref{prop:T-Brannen-Z3} wynika natychmiast $Q_K = 3/2$.
\hfill$\square$
\end{proof}
\end{theorem}

\subsection{Mechanizm fizyczny: dlaczego rozk\l{}ad r\'ownomierny?}

\begin{lemma}[Energia krzy\.zowa jest sta\l{}a na okr\k{e}gu $\mathcal{C}$]%
\label{lem:R2-Ecross-const}
Dla dowolnego rozmieszczenia
$\boldsymbol{\lambda} = (\lambda_1,\lambda_2,\lambda_3)$
na okr\k{e}gu~$\mathcal{C}$ (sta\l{}e $S_1$, $S_2$),
energia krzy\.zowa spe\l{}nia:
\begin{equation}
  E_{\rm cross} \;=\; \sum_{n<m}\lambda_n\lambda_m
  \;=\; \frac{S_1^2 - S_2}{2}
  \;=\; \mathrm{const}.
\end{equation}
\emph{Wniosek:} Minimalizacja $E_{\rm cross}$ \textbf{nie}
rozr\'o\.znia rozk\l{}ad\'ow na $\mathcal{C}$.
Potrzebny jest inny mechanizm selekcji
(skrypt \texttt{ex145}, test C8).
\end{lemma}

\begin{proposition}[Dekoherencja fazowa w~przestrzeni amplitud
$\Rightarrow$ $\mathbb{Z}_3$]%
\label{prop:R2-decoherence-Z3}
Niech $\boldsymbol{\lambda}$ le\.zy na okr\k{e}gu $\mathcal{C}$
i~niech $\alpha_k$ b\k{e}dzie k\k{a}tem odpowiadaj\k{a}cym
generacji $k$ na tym okr\k{e}gu (def.~\ref{def:R2-amplitude-vector}).

Koherencja fazowa sygna\l{}\'ow generacyjnych:
\begin{equation}\label{eq:R2-coherence}
  \mathcal{K}(\boldsymbol{\alpha})
  \;=\; \left|\sum_{k=0}^{2} e^{i\alpha_k}\right|^2
  \;=\; 3 + 2\sum_{n<m}\cos(\alpha_n - \alpha_m)
  \;\geq\; 0.
\end{equation}

\textbf{Zasada selekcji:} Maksymalna niezale\.zno\'s\'c
sektor\'ow generacyjnych (dekoherencja) wymaga
$\mathcal{K} = 0$, co jest r\'ownowa\.zne:
\begin{equation}\label{eq:R2-decoherence-condition}
  \sum_{n<m}\cos(\alpha_n - \alpha_m) = -\frac{3}{2}
  \qquad\text{(minimum globalne).}
\end{equation}

\medskip\noindent
\textbf{Jedyne rozwi\k{a}zanie:}
$\mathcal{K} = 0$ zachodzi wtedy i~tylko wtedy, gdy
k\k{a}ty tworz\k{a} r\'ownomierny rozk\l{}ad:
\begin{equation}\label{eq:R2-Z3-phases}
  \alpha_k = \alpha_0 + \frac{2\pi k}{3},
  \quad k = 0, 1, 2
\end{equation}
(modulo permutacje i~rotacj\k{e} globaln\k{a}).
Wynika to z~faktu, \.ze $\sum_k e^{i\alpha_k} = 0$
dla trzech wersork\'ow zachodzi \emph{wy\l{}\k{a}cznie}
przy symetrii $\mathbb{Z}_3$.

\begin{proof}
$|\sum_k e^{i\alpha_k}|^2 = 0$
$\Leftrightarrow$
$\sum_k e^{i\alpha_k} = 0$
$\Leftrightarrow$
trzy wektory jednostkowe $e^{i\alpha_k}$ na okr\k{e}gu
jednostkowym sumuj\k{a} si\k{e} do zera.
Dla $N=3$ jest to mo\.zliwe wtedy i~tylko wtedy,
gdy tworz\k{a} wierzcho\l{}ki tr\'ojk\k{a}ta r\'owno\-bocznego
--- tj.~$\alpha_k = \alpha_0 + 2\pi k/3$.
Nier\'owno\'s\'c $\mathcal{K} \geq 0$
jest nier\'owno\'sci\k{a} Cauchy'ego--Schwarza;
r\'owno\'s\'c $\mathcal{K} = 0$ to jedyne
globalne minimum.
\hfill$\square$
\end{proof}
\end{proposition}

\begin{remark}[Korekta wzgl\k{e}dem surowych faz ogonowych]%
\label{rem:R2-phase-correction}
Skrypt \texttt{ex125} pokaza\l{}, \.ze fazy ogonowe
$\delta_n$ profili $g^{(n)}(r)$ \textbf{nie} tworz\k{a}
post\k{e}pu $\mathbb{Z}_3$
($\delta_e = -89{,}2°$, $\delta_\mu = +78{,}1°$,
$\delta_\tau = +15{,}2°$).

Symetria $\mathbb{Z}_3$ działa na \textbf{k\k{a}tach
$\alpha_k$ w~przestrzeni amplitudowej}
$\boldsymbol{\lambda} = (A_0^2, A_1^2, A_2^2)$,
nie na surowych fazach $\delta_n$.
Weryfikacja: $\mathrm{CV}(A_n^2) = 1{,}000$
(ex126: 14/14~PASS) potwierdza r\'ownowa\.zno\'s\'c
z~symetri\k{a} $\mathbb{Z}_3$ w~amplitudach.
\end{remark}

\begin{remark}[Status T-OP1 po mo\'scie B2]
Ca\l{}y \l{}a\'ncuch wygl\k{a}da teraz:
\begin{enumerate}[nosep]
  \item Substrat $\Gamma$ z~$\mathbb{Z}_2$ $\Rightarrow$
    r\'ownanie solitonu z~potencja\l{}em $V(g)$.
  \item Trzy mody WKB ($n = 0, 1, 2$) $\Rightarrow$
    trzy generacje (prop.~\ref{prop:no-4th-generation}).
  \item Amplitudy ogonowe $A_n$ le\.z\k{a} na okr\k{e}gu
    $\mathcal{C}$ w~$\mathbb{R}^3$ (wi\k{e}z $S_1$, $S_2$).
  \item $E_{\rm cross} = \mathrm{const}$ na $\mathcal{C}$
    (lem.~\ref{lem:R2-Ecross-const}), wi\k{e}c selekcja
    wymaga kryterium fazowego.
  \item \textbf{Dekoherencja fazowa}
    $\mathcal{K} = |\sum_k e^{i\alpha_k}|^2 = 0$
    $\Rightarrow$ symetria $\mathbb{Z}_3$
    (prop.~\ref{prop:R2-decoherence-Z3}).
  \item Symetria $\mathbb{Z}_3$ $\Rightarrow$
    parametryzacja Brannena $\Rightarrow$ $Q_K = 3/2$
    (prop.~\ref{prop:T-Brannen-Z3}).
\end{enumerate}

\textbf{Status T-OP1}: \textbf{[Twierdzenie AN+NUM]} ---
\l{}a\'ncuch zamkni\k{e}ty z~zasad\k{a} dekoherencji fazowej
w~przestrzeni amplitud (nie surowych faz ogonowych).
Weryfikacja numeryczna: ex140 (12/12), ex126 (14/14),
ex145 (14/14).
\end{remark}

\begin{remark}[Weryfikacja numeryczna]\label{rem:R2-numerical}
Skrypt \texttt{ex140\_z3\_entropy\_optimum.py} weryfikuje:
\begin{enumerate}[nosep]
  \item Obliczenie faz ogonowych $\phi_n$ z~pe\l{}nego
    ODE solitonu (por.~ex125).
  \item Sprawdzenie warunku~\eqref{eq:R2-Z3-phases}.
  \item Obliczenie $Q_K$ z~amplitud numerycznych.
  \item Skan $\phi_0$ i~weryfikacja minimum $|E_{\rm int}|$.
\end{enumerate}

\textbf{Uwaga o bada\-niu ex125:}
Ex125 pokaza\l{}, \.ze fazy ogonowe ($\delta_e = -89{,}2°$,
$\delta_\mu = +78{,}1°$, $\delta_\tau = +15{,}2°$)
\textbf{nie} tworz\k{a} prostego post\k{e}pu $\mathbb{Z}_3$.
Symetria $\mathbb{Z}_3$ dzia\l{}a na k\k{a}tach $\alpha_k$
w~przestrzeni amplitudowej
(por.~rem.~\ref{rem:R2-phase-correction}).
\end{remark}

\subsection{Podsumowanie: zamkni\k{e}cie T-OP1}

\begin{center}
\begin{tabular}{@{}lll@{}}
\toprule
\textbf{Element} & \textbf{Status przed B2} & \textbf{Status po B2} \\
\midrule
$Q_K = 3/2$ algebraicznie & \statuslabel{Twierdzenie} &
  \statuslabel{Twierdzenie} \\
$\mathbb{Z}_3 \Rightarrow Q_K = 3/2$ (Brannen) &
  \statuslabel{Twierdzenie} & \statuslabel{Twierdzenie} \\
Dlaczego $\mathbb{Z}_3$ w~TGP? & \statuslabel{Program} &
  \statuslabel{Twierdzenie} \\
Pe\l{}ny \l{}a\'ncuch $\Gamma \to Q_K$ & \statuslabel{Otwarty} &
  \statuslabel{Twierdzenie} \\
\bottomrule
\end{tabular}
\end{center}
